\section{Documentary questions still under investigation}

This part records uncertainties; it does not weaken the hypotheses of
the checked theorems or certify an absent proof. The response to the
entire research prompt is still marked \textbf{INCOMPLETE} while these
questions and the last source audits are being pursued.

\subsection{The date and printed form of the Hilbert-space argument}

The inspected fourth printing of \emph{Produits tensoriels topologiques
et espaces nucl\'eaires} has a 1955 copyright and a 1966 printing date.
Chapter II, \S1, no.~4, printed p.~13, states the determinant expansion,
the eigenvalue-zero correspondence, and the genus-zero assertion for
Hilbert-space Schatten exponent at most one. Their mathematical
implication is proved in the checked part. The 1956 erratum, p.~119,
explains inadequate references to material removed from the Fredholm
paper. Neither inspected page dates the crossed-out manuscript sheets.

The original Bourbaki expos\'e 91 supplies a distinct textual witness:
\S4(c), printed p.~383, expressly states the Hilbert product formula
and coefficient identities, and the final page dates the text to
June 1957. Proposition~\ref{geo-H-history} gives the exact implication
for the trace formula. We do not identify that author date with the
date of an uninspected physical edition or with the date of another
manuscript.

The remaining dating question is kept separate from source qualifications. The
Montpellier catalogue describes the mathematical notes of cote~1 as
undated and identifies a letter dated 1953; the connection of that
letter with the crossed-out trace sheets must be inspected, not
assumed. The complete 150-page physical PDF of the dossier has now
been retrieved. Its cover is additional to the archive's 149 folios;
the physical PDF page number is normally one greater than the archive
folio number. Contact sheets of the whole dossier have been inspected,
but those thumbnails do not certify handwriting or dates.

The full memoir and the 1956 Fredholm article have now been checked
for the proof architecture, not merely for occurrences of a name.
The Hilbert genus-zero bridge is stated or used but not proved there;
the detailed source analysis is in Proposition~\ref{geo-H-history}.
The separate 1956 metric r\'esum\'e is receiving its final source audit.
The inspected memoir is the fourth printing, and line-for-line identity
with the first printing has not been established by collation. This
is an edition qualification, not an additional mathematical task
invented from the prompt. The publication is identified by its 1955
date while the inspected physical printing is identified as 1966.
Retrieval of Lidskii's original 1959 endnote is also under way.
The analytic proof in the checked part does not depend on resolving
the manuscript date. The cautious Lead that the sheets \emph{may}
predate 1959 is not replaced by a stronger priority assertion.

\subsection{What ``origin of SGA 6'' can mean}

The PS on 23/053--054, the final EGA IV 19.1.8, and SGA 6 VII
1.1--1.7 have been compared with original scans. The mathematical
correspondence is exact, and EGA's printed note refers to the 1966
seminar. The PS itself reports seminar work already under way and a
transitivity argument already checked. Consequently it does not
support a claim that the notion and proofs were first invented in
the act of writing that PS. A stronger claim about its causal
influence on the eventual SGA 6 wording would require a dated
document connecting the drafting stages; that fact is not inferred
from mathematical agreement alone.

\subsection{Scope distinctions, without additional conjectures}

The draft's earlier definition of condition P is absent from the retained
pages. We have therefore proved preservation of both explicitly defined
forms of formal projectivity in the noetherian adic Cohen-lifting theorem.
We do not identify an absent historical definition by assertion.
The natural-bundle theorem states its locality and finiteness hypotheses,
as the requested research package requires. Automatic formal locality
from mere \emph{\'etale} naturality is a stronger statement and is not
used. The programme headings in cote 26 are accompanied by precise
quantified formulations; no universal assertion about unspecified
``structures'' or unspecified moduli problems is invented.

All active states, failed source attempts, and next retrieval steps are
recorded in the dated ledger and the section-specific source logs.
